% GENERATED FILE. Edit canonical lessons, then run npm run generate:books.
% canonical-source-hash: fnv1a64:80bc536e2983ea5e
% canonical-lessons: LA-C192-cancer, LA-C192-passer, LA-C192-ciconia, LA-C192-noctua, LA-C192-pavo

\chapter{Crab, Sparrow, Stork, Owl, Peacock}
\label{ch:la-crab-sparrow-stork-owl-peacock}
\hlchaptermodality{192}

\begin{chapteropening}
\textbf{By the end of this chapter:} \emph{I can say a crab, a sparrow, a stork, an owl and a peacock.}

Complete the last lesson of chapter 192: I can say a crab, a sparrow, a stork, an owl and a peacock.
\end{chapteropening}

\section[cancer, cancrī]{cancer, cancrī — a crab}

\label{lesson:LA-C192-cancer}



\begin{quote}
Before the new one: say the Latin for a storm, a gale, then the Latin for a rain cloud.
\end{quote}

\subsection*{You'll want to know: cancer, cancrī}
\textbf{cancer, cancrī} — \textquotedblleft{}a crab\textquotedblright{}.

\begin{cousinweb}
\textbf{cancer} is probably inherited, from a doubled form of a root for \textquotedblleft{}hard\textquotedblright{}; Greek \emph{karkinos}, \textquotedblleft{}crab\textquotedblright{}, is a cousin (Lewis \& Short), and the n came in by dissimilation.
\end{cousinweb}

\begin{grammarlens}[title={What you've built}]
One more word for this chapter.
\end{grammarlens}

\subsection*{Your turn}
\begin{itemize}
\raggedright
  \item \emph{Say it:} \emph{cancer, cancrī}
  \item \emph{Say it:} \emph{cancer, cancrī}, once more
  \item \emph{Say it:} say \emph{cancer, cancrī}
  \item \emph{Recall:} say the Latin for a storm, a gale, then the Latin for a rain cloud, then say \emph{cancer, cancrī} again
\end{itemize}

\begin{tcolorbox}[breakable,colback=teal!4,colframe=teal!35!black,title={Before you move on}]
What does cancer, cancrī mean? (\textquotedblleft{}A crab\textquotedblright{}.) Say it once more.
\end{tcolorbox}
\section[passer, passeris]{passer, passeris — a sparrow}

\label{lesson:LA-C192-passer}



\begin{quote}
Before the new one: say the Latin for a rain cloud, then the Latin for a crab.
\end{quote}

\subsection*{You'll want to know: passer, passeris}
\textbf{passer, passeris} — \textquotedblleft{}a sparrow\textquotedblright{}.

\begin{cousinweb}
\textbf{passer} may be inherited, from a Proto-Indo-European \emph{*pettro-}, \textquotedblleft{}flier, bird\textquotedblright{}, from \emph{*peth2-}, \textquotedblleft{}to fly\textquotedblright{} (EtymDB 2.1); the details are uncertain.
\end{cousinweb}

\begin{grammarlens}[title={What you've built}]
One more word for this chapter.
\end{grammarlens}

\subsection*{Your turn}
\begin{itemize}
\raggedright
  \item \emph{Say it:} \emph{passer, passeris}
  \item \emph{Say it:} \emph{passer, passeris}, once more
  \item \emph{Say it:} say \emph{passer, passeris}
  \item \emph{Recall:} say the Latin for a rain cloud, then the Latin for a crab, then say \emph{passer, passeris} again
\end{itemize}

\begin{tcolorbox}[breakable,colback=teal!4,colframe=teal!35!black,title={Before you move on}]
What does passer, passeris mean? (\textquotedblleft{}A sparrow\textquotedblright{}.) Say it once more.
\end{tcolorbox}
\section[cicōnia, cicōniae]{cicōnia, cicōniae — a stork}

\label{lesson:LA-C192-ciconia}



\begin{quote}
Before the new one: say the Latin for a crab, then the Latin for a sparrow.
\end{quote}

\subsection*{You'll want to know: cicōnia, cicōniae}
\textbf{cicōnia, cicōniae} — \textquotedblleft{}a stork\textquotedblright{}.

\begin{cousinweb}
\textbf{cicōnia} is probably a doubled form of the root \emph{*keh2n-}, \textquotedblleft{}to sing, to cry\textquotedblright{} (EtymDB 2.1). At Praeneste the bird was called \emph{cōnia}, \textquotedblleft{}stork\textquotedblright{} (Lewis \& Short), and \emph{canō}, \textquotedblleft{}to sing\textquotedblright{}, has the same root.
\end{cousinweb}

\begin{grammarlens}[title={What you've built}]
One more word for this chapter.
\end{grammarlens}

\subsection*{Your turn}
\begin{itemize}
\raggedright
  \item \emph{Say it:} \emph{cicōnia, cicōniae}
  \item \emph{Say it:} \emph{cicōnia, cicōniae}, once more
  \item \emph{Say it:} say \emph{cicōnia, cicōniae}
  \item \emph{Recall:} say the Latin for a crab, then the Latin for a sparrow, then say \emph{cicōnia, cicōniae} again
\end{itemize}

\begin{tcolorbox}[breakable,colback=teal!4,colframe=teal!35!black,title={Before you move on}]
What does cicōnia, cicōniae mean? (\textquotedblleft{}A stork\textquotedblright{}.) Say it once more.
\end{tcolorbox}
\section[noctua, noctuae]{noctua, noctuae — an owl}

\label{lesson:LA-C192-noctua}



\begin{quote}
Before the new one: say the Latin for a sparrow, then the Latin for a stork.
\end{quote}

\subsection*{You'll want to know: noctua, noctuae}
\textbf{noctua, noctuae} — \textquotedblleft{}an owl\textquotedblright{}.

\begin{cousinweb}
\textbf{noctua} is built on \emph{nox}, \emph{noctis}, \textquotedblleft{}night\textquotedblright{} (Lewis \& Short): the night bird.
\end{cousinweb}

\begin{grammarlens}[title={What you've built}]
One more word for this chapter.
\end{grammarlens}

\subsection*{Your turn}
\begin{itemize}
\raggedright
  \item \emph{Say it:} \emph{noctua, noctuae}
  \item \emph{Say it:} \emph{noctua, noctuae}, once more
  \item \emph{Say it:} say \emph{noctua, noctuae}
  \item \emph{Recall:} say the Latin for a sparrow, then the Latin for a stork, then say \emph{noctua, noctuae} again
\end{itemize}

\begin{tcolorbox}[breakable,colback=teal!4,colframe=teal!35!black,title={Before you move on}]
What does noctua, noctuae mean? (\textquotedblleft{}An owl\textquotedblright{}.) Say it once more.
\end{tcolorbox}
\section[pāvō, pāvōnis]{pāvō, pāvōnis — a peacock}

\label{lesson:LA-C192-pavo}



\begin{quote}
Before the new one: say the Latin for a stork, then the Latin for an owl.
\end{quote}

\subsection*{You'll want to know: pāvō, pāvōnis}
\textbf{pāvō, pāvōnis} — \textquotedblleft{}a peacock\textquotedblright{}.

\begin{cousinweb}
\textbf{pāvō} is probably a loan from the east, like Greek \emph{taōs}, \textquotedblleft{}peacock\textquotedblright{} (Lewis \& Short); both are thought to go back to a language of India, but the path is uncertain. English \emph{peacock} begins with it.
\end{cousinweb}

\begin{grammarlens}[title={What you've built}]
One more word for this chapter.
\end{grammarlens}

\subsection*{Your turn}
\begin{itemize}
\raggedright
  \item \emph{Say it:} \emph{pāvō, pāvōnis}
  \item \emph{Say it:} \emph{pāvō, pāvōnis}, once more
  \item \emph{Say it:} say \emph{pāvō, pāvōnis}
  \item \emph{Recall:} say the Latin for a stork, then the Latin for an owl, then say \emph{pāvō, pāvōnis} again
\end{itemize}

\begin{tcolorbox}[breakable,colback=teal!4,colframe=teal!35!black,title={Before you move on}]
What does pāvō, pāvōnis mean? (\textquotedblleft{}A peacock\textquotedblright{}.) Say it once more.
\end{tcolorbox}
